% 数学建模竞赛论文 LaTeX 模板
% 使用 xelatex 编译: xelatex paper.tex

\documentclass[12pt,a4paper]{article}

% ---- 中文支持 ----
\usepackage{ctex}

% ---- 页面设置 ----
\usepackage[top=2.54cm, bottom=2.54cm, left=3.18cm, right=3.18cm]{geometry}
\usepackage{setspace}
\onehalfspacing

% ---- 数学公式 ----
\usepackage{amsmath, amssymb, amsthm}

% ---- 图表 ----
\usepackage{graphicx}
\usepackage{caption}
\usepackage{subcaption}
\usepackage{booktabs}  % 三线表

% ---- 引用 ----
\usepackage[numbers,sort&compress]{natbib}

% ---- 超链接 ----
\usepackage{hyperref}
\hypersetup{
    colorlinks=true,
    linkcolor=black,
    citecolor=black,
    urlcolor=blue
}

% ---- 编号 ----
\usepackage{chngcntr}
\counterwithin{figure}{section}
\counterwithin{table}{section}
\counterwithin{equation}{section}

% ---- 自定义命令 ----
\newcommand{\figref}[1]{图~\ref{#1}}
\newcommand{\tabref}[1]{表~\ref{#1}}
\newcommand{\eqref}[1]{式~(\ref{#1})}

% ============================================================
\begin{document}

% ---- 标题 ----
\begin{center}
    \LARGE\textbf{数学建模竞赛论文}\\
    \vspace{0.5cm}
    \large 参赛编号: XXXXX\\
    \vspace{0.5cm}
\end{center}

% ---- 摘要 ----
\begin{abstract}
\noindent 在此填写摘要内容。摘要应包含：问题概述、方法概述、主要结果。

\noindent\textbf{关键词}：关键词1；关键词2；关键词3
\end{abstract}

\thispagestyle{empty}
\newpage

% ---- 目录 ----
\tableofcontents
\newpage

% ============================================================
\section{问题重述与分析}
\label{sec:intro}

\subsection{问题重述}
\label{sec:restate}

在此重述题目要求\cite{ref1}。

\subsection{问题分析}
\label{sec:analysis}

在此分析解题思路。

% ============================================================
\section{模型假设}
\label{sec:assumptions}

\begin{enumerate}
    \item 假设一：\ldots
    \item 假设二：\ldots
    \item 假设三：\ldots
\end{enumerate}

% ============================================================
\section{符号说明}
\label{sec:notation}

\begin{table}[htbp]
    \centering
    \caption{主要符号说明}
    \begin{tabular}{ccc}
        \toprule
        符号 & 含义 & 单位 \\
        \midrule
        $x$ & \ldots & - \\
        $y$ & \ldots & - \\
        \bottomrule
    \end{tabular}
\end{table}

% ============================================================
\section{模型建立与求解}
\label{sec:model}

\subsection{问题一模型}
\label{sec:prob1}

在这里建立数学模型：

\begin{equation}
    \label{eq:model}
    \min f(x) = \sum_{i=1}^{n} c_i x_i
\end{equation}

约束条件：
\begin{align}
    \sum_{i=1}^{n} a_{ij} x_i &\leq b_j, \quad j = 1,\ldots,m \\
    x_i &\geq 0, \quad i = 1,\ldots,n
\end{align}

\begin{figure}[htbp]
    \centering
    \includegraphics[width=0.7\textwidth]{figures/fig1_功率曲线.png}
    \caption{功率曲线}
    \label{fig:power}
\end{figure}

\subsection{问题二模型}
\label{sec:prob2}

\ldots

% ============================================================
\section{模型检验与灵敏度分析}
\label{sec:validation}

\subsection{模型检验}
\label{sec:verify}

\subsection{灵敏度分析}
\label{sec:sensitivity}

% ============================================================
\section{模型的优缺点与改进方向}
\label{sec:discussion}

\subsection{模型优点}
\begin{itemize}
    \item 优点一
    \item 优点二
\end{itemize}

\subsection{模型缺点}
\begin{itemize}
    \item 缺点一
    \item 缺点二
\end{itemize}

\subsection{改进方向}
\begin{itemize}
    \item 改进方向一
    \item 改进方向二
\end{itemize}

% ============================================================
\begin{thebibliography}{99}
\bibitem{ref1} 作者. 标题[J]. 期刊, 年份, 卷(期): 页码.
\end{thebibliography}

\end{document}
